\section{Comparison to the Dark Energy Survey}\label{sec:des}

The Dark Energy Survey (DES) was a six-year optical and near-infrared imaging survey designed to investigate the origin of cosmic acceleration. 
DES operated from 2013 to 2019 using the Dark Energy Camera (DECam) mounted on the \SI{4}{\m} Victor M.\ Blanco Telescope at Cerro Tololo Inter-American Observatory in Chile. 
DECam, the camera used to conduct the survey, is a 570-megapixel wide-field imager with a $\sim3\,\mathrm{deg}^2$ field of view, consisting of 62 science CCDs. 
The survey imaged approximately \SI{5000}{deg^2} of the southern sky in five broadband filters ({\it grizY}) to a typical depth of $i \approx 24$\,mag (10$\sigma$ for galaxies). 

The PSF pipelines used by the Rubin Observatory are modelled after the DES pipelines, including the use of the same PSF modelling software: \texttt{PiFF}.
As the two surveys share similar science goals, we compare current Rubin Observatory PSF characteristics to DES Year 6 PSF results \citep{schutt2025}. The DES Y6 PSF star sample consists of $\sim$ 110 million stars, while the Rubin week 49 DRP data contains $\sim$ 153 million source detections.
Note that the DES Y6 colour fitting was done with a first order polynomial, while this work uses a second order polynomial.

\begin{figure}[h!]
    \centering
\includegraphics[width=0.4\textwidth]{figures/des_mag.pdf}
    \caption{DES Y6 PSF $\delta T/T$, $\delta e_1$ and $\delta e_2$ as a function of magnitude. Stars are divided into equal-sized colour bins of $\textit{g}-\textit{i} \in [0.0, 0.8, 2.5, 3.5]$
    (This plot is identical to the top three panels of Fig.~15 in \cite{schutt2025}).} 
    \label{fig:des_mag}
\end{figure}

Figure~\ref{fig:des_mag} shows the dependence of chromatically-fitted PSF $\delta T/T$, $\delta e_1$ and $\delta e_2$ residuals on magnitude in DES Y6 data. 
Comparing the residuals to Fig.~\ref{subfig:residuals_vs_mag_b}, we find that the magnitude of \texttt{PiFF} residuals are on par with DES levels, even though the overall trend looks different.

\begin{figure}[h!]
    \centering
\includegraphics[width=0.5\textwidth]{figures/des_color.pdf}
    \caption{DES Y6 residuals as a function of $\textit{g}-\textit{i}$ colour (for the {\it gri} bands) and $\textit{i}-\textit{z}$ colour (for the {\it z} band)
    (This plot is identical to the four left panels of Fig.~18 in \cite{schutt2025}).} 
    \label{fig:des_color}
\end{figure}

Figure~\ref{fig:des_color} shows the DES Y6 PSF residuals as a function of $\textit{g}-\textit{i}$ and $\textit{i}-\textit{z}$ colour. 
Comparing the colour dependence to Fig.~\ref{subfig:color_vs_residuals_b}, DES Y6 {\it g} band $\delta T/T$ residuals are up to $\sim$ 50\% larger than Rubin.

\begin{figure}[h!]
  \centering

  \begin{subfigure}[t]{0.48\textwidth}
    \centering
    \includegraphics[width=\linewidth]{figures/des_dcr_g.pdf}
    \label{subfig:des_dcr_g}
  \end{subfigure}
  \hfill
  \begin{subfigure}[t]{0.48\textwidth}
    \centering
    \includegraphics[width=\linewidth]{figures/des_dcr_r.pdf}
    \label{subfig:des_dcr_r}
  \end{subfigure}

  \caption{DES Y6 PSF $\delta e_1$ and $\delta e_2$ as a function of $\mathrm{DCR_1}$ and $\mathrm{DCR_2}$ respectively, in the {\it g} and {\it r} bands
  (This plot is identical to the top two panels of Fig.~16 in \cite{schutt2025}).}
  \label{fig:des_dcr}
\end{figure}

Figure~\ref{fig:des_dcr} shows the DES Y6 ellipticity residuals as a function of $\mathrm{DCR_1}$ and $\mathrm{DCR_2}$ in the {\it g} and {\it r} bands. 
Comparing these to Fig.~\ref{subfig:g_dcr_b} and Fig.~\ref{subfig:r_dcr_b}, we find that in the Rubin data, the PSF model errors coming from DCR are $\sim$ 80\% lower than that in DES Y6 in the {\it g} band and to similar levels in the {\it r} band.

